% !TEX program = platex
% 数学ノート集・共通テンプレート v2（積分論2026年度第1部の版面を参考）
% A5判。章内で定義・定理・例・注意等の番号を共通化。
% コンパイル： platex -interaction=nonstopmode math_note_template_v2.tex (2回)
%              dvipdfmx math_note_template_v2.dvi
\documentclass[a5paper,10pt,onecolumn]{jsarticle}
\usepackage[dvipdfmx,top=22mm,bottom=20mm,inner=19mm,outer=18mm,headheight=13pt,headsep=6mm]{geometry}
\usepackage{amsmath,amssymb,amsthm,mathtools}
\usepackage{enumitem,fancyhdr,titlesec}
\usepackage[dvipdfmx,bookmarks=true,colorlinks=false]{hyperref}
\usepackage{pxjahyper}

% ----- 差し替える文書情報 -----
\newcommand{\NoteTitle}{凹角領域における楕円型正則性の限界}
\newcommand{\NoteSeries}{数学ノート集}

% ----- 文字・組版 -----
% 参考PDFの質素な明朝＋Computer Modernの印象を引き継ぎながら，
% A5判で窮屈にならないよう本文の行間を少し確保する．
\linespread{1.13}
\setlength{\parindent}{1zw}
\setlength{\parskip}{0pt}
\setlength{\abovedisplayskip}{6pt plus 2pt minus 2pt}
\setlength{\belowdisplayskip}{6pt plus 2pt minus 2pt}
\setlength{\abovedisplayshortskip}{3pt plus 2pt minus 1pt}
\setlength{\belowdisplayshortskip}{4pt plus 2pt minus 1pt}
\setlist[enumerate]{topsep=3pt,itemsep=1pt,parsep=0pt,leftmargin=2.4zw}
\titleformat{\section}{\normalfont\normalsize\bfseries}{\thesection}{0.8zw}{}
\titlespacing*{\section}{0pt}{2.0ex plus .4ex}{.9ex}
\titleformat{\subsection}{\normalfont\normalsize\bfseries}{\thesubsection}{0.7zw}{}
\titlespacing*{\subsection}{0pt}{1.5ex plus .3ex}{.6ex}
\raggedbottom

% ----- 柱とページ番号：ページ上部、罫線なし -----
\fancypagestyle{mathnote}{%
  \fancyhf{}%
  \fancyhead[L]{\footnotesize \NoteSeries\ （\NoteTitle）}%
  \fancyhead[R]{\footnotesize \thepage}%
  \renewcommand{\headrulewidth}{0pt}%
  \renewcommand{\footrulewidth}{0pt}%
}
\pagestyle{mathnote}

% ----- 節内通し番号の定義・定理群 -----
\newtheoremstyle{mathnotetheorem}
 {4.5pt}{4.5pt}{\normalfont}{}%
 {\bfseries}{}{0.65zw}%
 {\thmname{#1}\thmnumber{ #2}\thmnote{（#3）}}
\theoremstyle{mathnotetheorem}
\newtheorem{theorem}{定理}[section]
\newtheorem{definition}[theorem]{定義}
\newtheorem{lemma}[theorem]{補題}
\newtheorem{proposition}[theorem]{命題}
\newtheorem{corollary}[theorem]{系}
\newtheorem{example}[theorem]{例}
\newtheorem{remark}[theorem]{注意}
\renewcommand{\proofname}{証明}
\renewcommand{\qedsymbol}{$\square$}
\makeatletter
\renewenvironment{proof}[1][\proofname]{%
  \par\pushQED{\qed}\normalfont\topsep4\p@\@plus2\p@\relax
  \trivlist\item[\hskip\labelsep\bfseries #1]\ignorespaces
}{\popQED\endtrivlist\@endpefalse}
\makeatother

% ----- 使うノートで適宜拡張する記号 -----
\newcommand{\R}{\mathbb{R}}
\newcommand{\norm}[1]{\lVert #1\rVert}
\newcommand{\abs}[1]{\lvert #1\rvert}

\begin{document}
\section*{読了後にできること}
本稿を読むための前提は，Sobolev空間，極座標，楕円型境界値問題．
本稿の具体的な到達課題は，凹角の特異関数が $H^1$ だが $H^2$ でないことを積分で判定する．

\section{問題と主定理}
本稿では，境界のLipschitz性だけからPoisson方程式の弱解の
$H^2$ 正則性は従わないことを，具体的な函数と右辺を構成して証明する．
この事実は，有限要素法の誤差評価で厳密解の$H^2$ 正則性を
独立した仮定として置く必要があることを示す．

\begin{definition}[弱解]
有界開集合$\Omega\subset\mathbb R^2$と$f\in L^2(\Omega)$に対し，
$u\in H_0^1(\Omega)$が
\[
 \int_\Omega\nabla u\cdot\nabla v\,dx
   =\int_\Omega fv\,dx
 \qquad(v\in H_0^1(\Omega))
\]
を満たすとき，$-\Delta u=f$の零Dirichlet境界条件に対する弱解という．
ここで$H_0^1(\Omega)$は$C_c^\infty(\Omega)$の$H^1$閉包である．
\end{definition}

\begin{theorem}[凹角領域における$H^2$正則性の破綻]\label{thm:corner}
$\pi<\omega<2\pi$とし，
\[
 \Omega_\omega=
 \{(r\cos\theta,r\sin\theta):0<r<1,\ 0<\theta<\omega\}
\]
と定める．$\Omega_\omega$は有界Lipschitz領域である．
ある$f\in C^\infty(\overline{\Omega_\omega})$が存在し，
零Dirichlet境界条件を持つ$-\Delta u=f$の弱解は
ただ一つ存在するが，その解は$H^2(\Omega_\omega)$に属さない．
$f$は原点の近傍で恒等的に零となるように選べる．
\end{theorem}

この定理における$C^\infty(\overline\Omega)$は，
閉包の各点の近傍で滑らかな函数へ延長できることを意味する．

\section{領域の形状と角点函数}
$\lambda=\pi/\omega$とおく．$1/2<\lambda<1$である．

\begin{lemma}[扇形のLipschitz境界]\label{lem:lipschitz}
$\Omega_\omega$は有界Lipschitz領域である．
\end{lemma}
\begin{proof}
原点以外の境界は，二本の直線の一部と半径1の円弧からなり，
それらの交点では二つの滑らかな曲線が零でない角度で交わる．
従って各点の近傍で適切な直交座標を取れば，境界はLipschitz函数の
グラフで，領域はその一方の側として表される．
原点について確認する．領域の補集合が占める角度を
$\beta=2\pi-\omega\in(0,\pi)$とする．欠損角の二等分線を
負の$t$軸とする直交座標$(s,t)$を取る．十分小さい原点近傍では，
領域は
\[
 t>-\cot(\beta/2)|s|
\]
で表される．右辺はLipschitz函数である．
境界のコンパクト性から有限個のグラフ近傍で覆える．
\end{proof}

\begin{lemma}[角点の調和函数]\label{lem:harmonic}
$0<r<1$, $0<\theta<\omega$で
\[
 w(r,\theta)=r^\lambda\sin(\lambda\theta)
\]
と置くと，$w$は滑らかで$\Delta w=0$となる．さらに
$w(r,0)=w(r,\omega)=0$である．
\end{lemma}
\begin{proof}
極座標におけるLaplacianは
\[
 \Delta=\partial_{rr}+r^{-1}\partial_r+r^{-2}\partial_{\theta\theta}
\]
である．これに$w$を代入すると
\[
 \Delta w=(\lambda(\lambda-1)+\lambda-\lambda^2)
 r^{\lambda-2}\sin(\lambda\theta)=0.
\]
境界値については$\sin0=\sin(\lambda\omega)=\sin\pi=0$
を用いればよい．
\end{proof}

\section{弱解の構成と零境界条件}
$\chi\in C^\infty([0,\infty))$を
\[
 0\le\chi\le1,\quad
 \chi(r)=1\ (0\le r\le1/3),\qquad
 \chi(r)=0\ (r\ge2/3)
\]
となるように取る．以後
\[
 u(r,\theta)=\chi(r)r^\lambda\sin(\lambda\theta)
\]
と定める．

\begin{lemma}[エネルギー有限性]\label{lem:H1}
$u\in H^1(\Omega_\omega)$であり，原点近傍では
\[
 \int_{\Omega_\omega\cap\{r<\varepsilon\}}|u|^2\,dx
      \le C\varepsilon^{2\lambda+2},\qquad
 \int_{\Omega_\omega\cap\{r<\varepsilon\}}|\nabla u|^2\,dx
      \le C\varepsilon^{2\lambda}
\]
が$0<\varepsilon<1/3$について成立する．
\end{lemma}
\begin{proof}
$r<1/3$では
\[
 |u|\le r^\lambda,\qquad
 |\nabla u|^2=(\partial_r u)^2+r^{-2}(\partial_\theta u)^2
             =\lambda^2r^{2\lambda-2}.
\]
極座標の面積要素$r\,dr\,d\theta$を用いて積分すると
\[
 \int_{r<\varepsilon}|u|^2\,dx
   \le\omega\int_0^\varepsilon r^{2\lambda+1}\,dr,
 \qquad
 \int_{r<\varepsilon}|\nabla u|^2\,dx
   =\omega\lambda^2\int_0^\varepsilon r^{2\lambda-1}\,dr.
\]
$\lambda>0$なので両積分は有限であり，記載した評価を得る．
$r\ge1/3$で$u$とその導函数は有界で，$r\ge2/3$では零である．
原点は領域の境界点であるから，領域内部での古典的な微分は
弱微分と一致し，$u\in H^1(\Omega_\omega)$を得る．
\end{proof}

\begin{lemma}[零境界条件]\label{lem:H0}
$u\in H_0^1(\Omega_\omega)$である．
\end{lemma}
\begin{proof}
$\eta\in C^\infty([0,\infty))$を
$0\le\eta\le1$，$\eta(t)=0$（$t\le1$），
$\eta(t)=1$（$t\ge2$）となるように取る．
$u_\varepsilon=\eta(r/\varepsilon)u$と置く．
前補題の評価と
$|\nabla\eta(r/\varepsilon)|\le C/\varepsilon$から
\[
 \|u-u_\varepsilon\|_{H^1(\Omega_\omega)}
       \longrightarrow0\qquad(\varepsilon\downarrow0)
\]
を得る．実際，$\nabla(u-u_\varepsilon)$の第一項は
$r<2\varepsilon$上の$\nabla u$で抑えられ，第二項は
$\varepsilon<r<2\varepsilon$上で
$C\varepsilon^{-1}|u|\le C\varepsilon^{\lambda-1}$
と抑えられるので，それぞれの$L^2$ノルムは
$O(\varepsilon^\lambda)$である．

$\varepsilon>0$を固定する．滑らかな角度切断函数
$\zeta_\delta(\theta)$を
$\theta\in(0,\delta)\cup(\omega-\delta,\omega)$では0，
$\theta\in(2\delta,\omega-2\delta)$では1，
$|\zeta_\delta'|\le C/\delta$となるように取る．
$u_{\varepsilon,\delta}=\zeta_\delta(\theta)u_\varepsilon$
は，$0<\delta<\omega/4$について$C_c^\infty(\Omega_\omega)$に属する．
$u_\varepsilon$の台は$r\ge\varepsilon$にあり，
$\sin(\lambda\theta)$は両端角の近くで
$C\min(\theta,\omega-\theta)$以下である．
従って$\zeta_\delta'$を含む項
$r^{-1}\zeta_\delta' u_\varepsilon$の絶対値は
$r\ge\varepsilon$上で$C_\varepsilon$以下である．
その台の角度方向の幅は$O(\delta)$であるため，
この項の$L^2$ノルムは$O_\varepsilon(\sqrt\delta)$となる．
他の項は有界収束定理で収束する．ゆえに
\[
 \|u_{\varepsilon,\delta}-u_\varepsilon\|_{H^1}
   \longrightarrow0\qquad(\delta\downarrow0).
\]
対角列を選べば$u$を$C_c^\infty(\Omega_\omega)$で
$H^1$近似できる．従って$u\in H_0^1(\Omega_\omega)$である．
\end{proof}

\section{滑らかな右辺と弱解の一意性}
\begin{lemma}[右辺の正則性]\label{lem:rhs}
$f=-\Delta u$を$\Omega_\omega$上で古典的に定めると，
$f$は閉包の近傍に滑らかに延長でき，
原点近傍で零となる．さらに
\[
 f(r,\theta)=-\bigl(\chi''(r)r^\lambda
      +(2\lambda+1)\chi'(r)r^{\lambda-1}\bigr)
      \sin(\lambda\theta)
\]
が成り立つ．
\end{lemma}
\begin{proof}
積の微分公式と補題\ref{lem:harmonic}から
\[
 \Delta(\chi w)=(\chi''+r^{-1}\chi')w
                +2\chi'\partial_rw
\]
となる．$\partial_rw=\lambda r^{\lambda-1}\sin(\lambda\theta)$を
代入して表示式を得る．$\chi',\chi''$は
$1/3\le r\le2/3$に台を持つため，$f$も同じ帯に台を持つ．
この帯では極座標表示が境界上を含めて滑らかであり，
$f$は$C^\infty(\overline{\Omega_\omega})$に属する．
\end{proof}

\begin{lemma}[弱形式と一意性]\label{lem:weak}
上で構成した$u$は，$f=-\Delta u$の零Dirichlet弱解であり，
そのような弱解はただ一つである．
\end{lemma}
\begin{proof}
$u$は領域内部で滑らかで，$f=-\Delta u$が成立する．
任意の$v\in C_c^\infty(\Omega_\omega)$に対し
通常の部分積分で
\[
 \int_{\Omega_\omega}\nabla u\cdot\nabla v\,dx
   =\int_{\Omega_\omega}fv\,dx
\]
を得る．左辺は$H_0^1$，右辺は$L^2$ノルムに関して連続である．
$C_c^\infty(\Omega_\omega)$の$H_0^1$稠密性から
全ての$v\in H_0^1(\Omega_\omega)$へ等式が拡張される．

一意性を示す．同じ$f$に対する弱解$u_1,u_2$を取り，
$w=u_1-u_2$と置くと，$v=w$を代入して
$\|\nabla w\|_2^2=0$を得る．
$H_0^1(\Omega_\omega)$上のPoincar\'e不等式により
$\|w\|_2=0$となる．この不等式は領域を含む長方形で
$C_c^\infty(\Omega_\omega)$を零延長し，一変数の
Newton--Leibniz公式とCauchy--Schwarzを適用して導ける．
従って$u_1=u_2$である．
\end{proof}

\section{二階微分の非可積分性}
\begin{lemma}[角点特異性の強さ]\label{lem:notH2}
構成した$u$は$H^2(\Omega_\omega)$に属さない．
\end{lemma}
\begin{proof}
$r<1/3$では$\chi=1$であり，
\[
 \partial_{rr}u
   =\lambda(\lambda-1)r^{\lambda-2}\sin(\lambda\theta).
\]
$\lambda\in(1/2,1)$なので$\lambda(\lambda-1)\ne0$である．
$\omega/4\le\theta\le3\omega/4$では
$\sin(\lambda\theta)\ge1/\sqrt2$である．従って
\[
 \int_{\Omega_\omega\cap\{r<1/3\}}
       |\partial_{rr}u|^2\,dx
 \ge C\int_0^{1/3}r^{2\lambda-3}\,dr=+\infty,
\]
最後の発散は$\lambda<1$から従う．

もし$u\in H^2(\Omega_\omega)$なら，領域内部での古典的な
二階微分は弱二階微分と一致し，
\[
 \partial_{rr}u
 =\cos^2\theta\,\partial_{xx}u
  +2\sin\theta\cos\theta\,\partial_{xy}u
  +\sin^2\theta\,\partial_{yy}u
\]
であるから$\partial_{rr}u\in L^2(\Omega_\omega)$となる．
これは前の発散に矛盾する．
\end{proof}

\begin{proof}[定理\ref{thm:corner}の証明]
補題\ref{lem:lipschitz}によって領域は有界Lipschitzである．
補題\ref{lem:H0}により構成した$u$は零境界条件を持つ．
補題\ref{lem:rhs}により$f$は滑らかであり，
補題\ref{lem:weak}から弱解の存在と一意性を得る．
補題\ref{lem:notH2}からその解は$H^2$に属さない．
\end{proof}

\section{正則性定理と有限要素法への帰結}
\begin{corollary}[Lipschitz境界だけでは不十分]
すべての有界Lipschitz領域$\Omega$と全ての
$f\in L^2(\Omega)$について，零Dirichlet問題の弱解が
$H^2(\Omega)$に属するという命題は偽である．
また，一般のLipschitz領域で
\[
 \|u\|_{H^2(\Omega)}\le C_\Omega\|f\|_{L^2(\Omega)}
\]
というアプリオリ評価を主張することもできない．
\end{corollary}
\begin{proof}
定理\ref{thm:corner}の領域と滑らかな$f$を選べば，
左辺は有限値ではないので両命題の反例となる．
\end{proof}

\begin{remark}[有限要素法への適用]
形状正則三角形分割に対する一次有限要素補間評価
\[
 \|u-I_hu\|_{H^1(\Omega)}\le Ch|u|_{H^2(\Omega)}
\]
は，$u\in H^2(\Omega)$という仮定の下で成立する評価である．
本稿の$u$には右辺が有限でないため，この評価を適用できない．
ここから直ちに有限要素解が収束しないという結論は出ない．
より弱い正則性を仮定した誤差評価は別途検討を要する．
\end{remark}

\begin{remark}[論理的な結論]
Poisson方程式の$H_0^1$弱解の存在一意性と$H^2$正則性は
異なる主張である．前者はこの反例上でも成立するが後者は失敗する．
「一般Lipschitz領域での$H^2$正則性」を完成させる代わりに，
正則性定理には境界形状，係数および右辺について
実際に十分な仮定を付ける必要がある．
\end{remark}

\begin{thebibliography}{9}
\bibitem{grisvard} P. Grisvard,
 \emph{Elliptic Problems in Nonsmooth Domains}, Pitman, 1985.
\bibitem{evans} L. C. Evans,
 \emph{Partial Differential Equations}, 2nd ed., AMS, 2010.
\end{thebibliography}

\section{到達課題の完全解答}
\noindent\textbf{凹角で二階正則性が失われる理由．}
\begin{proposition}
開扇形
$\Omega=\{(r,\theta):0<r<1,\ 0<\theta<\omega\}$
を $\pi<\omega<2\pi$ として考え，
$\alpha=\pi/\omega\in(1/2,1)$ と置く．
\[
u(r,\theta)=r^\alpha\sin(\alpha\theta)
\]
は原点近傍で $H^1$ に属するが $H^2$ には属さない．
\end{proposition}
\begin{proof}
極座標の微分から
$|\nabla u|^2=\alpha^2r^{2\alpha-2}$．
原点近くのエネルギー積分は
定数倍の $\int_0^\varepsilon r^{2\alpha-1}dr$
となり，$\alpha>0$ だから有限である．
一方二階微分は一般に
$r^{\alpha-2}$ の大きさを持つ．
その二乗の面積積分は
$\int_0^\varepsilon r^{2\alpha-3}dr$
を含む．$\alpha<1$ より指数が $-1$ 未満となり発散する．
\end{proof}
\begin{remark}
原点から離れた場所で滑らかな径方向のカットオフを掛ければ，
零境界値を持つ $H_0^1$ 関数で，
Laplacianが $L^2$ に属するにもかかわらず
$H^2$ に属さない例を得る．
\end{remark}
\noindent\textbf{到達確認．}
角度 $\omega=3\pi/2$ のとき $\alpha$ はいくつか．
\par\noindent\textbf{解答．}
$\alpha=\pi/(3\pi/2)=2/3$ であり，
$2\alpha-3=-5/3<-1$ だから二階微分の二乗は積分不能である．
\end{document}